% Ficha de estudio 4 del Taller virtual ATA EPE: ligaduras dinámicas, tornillos, tuercas y cuerpos compuestos.
\documentclass[11pt,a4paper]{article}
% Preámbulo común de las fichas de estudio del Taller virtual ATA EPE.
\usepackage[spanish,es-noshorthands,es-tabla]{babel}
\usepackage[utf8]{inputenc}
\usepackage[T1]{fontenc}
\usepackage{lmodern}
\usepackage[margin=2.2cm]{geometry}
\usepackage[table]{xcolor}
\usepackage{amsmath,amssymb,bm}
\usepackage{tabularx,array,booktabs}
\usepackage{enumitem}
\usepackage{fontawesome5}
\usepackage[most]{tcolorbox}
\usepackage[hidelinks]{hyperref}
\usepackage{listings}

\definecolor{morado}{HTML}{2A1B7E}
\definecolor{amarillo}{HTML}{FFCF2E}
\definecolor{rojo}{HTML}{C4142C}
\definecolor{verde}{HTML}{1F9D55}
\definecolor{gris}{HTML}{5D5A73}

\newcolumntype{Y}{>{\raggedright\arraybackslash}X}
\renewcommand{\arraystretch}{1.25}
\setlist{itemsep=1pt,topsep=3pt}
\numberwithin{equation}{section}
\lstset{basicstyle=\ttfamily\small,columns=fullflexible,keepspaces=true,frame=single,rulecolor=\color{gris!40},
  literate={á}{{\'a}}1 {é}{{\'e}}1 {í}{{\'i}}1 {ó}{{\'o}}1 {ú}{{\'u}}1 {ñ}{{\~n}}1 {Ω}{{$\Omega$}}1
    {θ}{{$\theta$}}1 {λ}{{$\lambda$}}1 {Δ}{{$\Delta$}}1 {μ}{{$\mu$}}1 {·}{{$\cdot$}}1 {←}{{$\leftarrow$}}1 {§}{{\S}}1 {→}{{$\to$}}1 {≠}{{$\neq$}}1 {É}{{\'E}}1 {Á}{{\'A}}1 {∨}{{$\vee$}}1 {∧}{{$\wedge$}}1 {σ}{{$\sigma$}}1 {×}{{$\times$}}1 {≈}{{$\approx$}}1}

\newcommand{\rev}[1]{\widetilde{#1}}
\newcommand{\I}{\mathbf{I}}
\newcommand{\vect}[1]{\bm{#1}}
\newcommand{\A}{\mathcal{A}}

\newtcolorbox{meta}{colback=morado,colframe=morado,coltext=white,arc=3mm}
\newtcolorbox{lectura}{colback=morado!5,colframe=morado,fonttitle=\bfseries,title={\faBook~Lecturas}}
\newtcolorbox{ejercicio}[1]{colback=amarillo!10,colframe=amarillo!80!black,fonttitle=\bfseries,title={\faPencil*~Ejercicio #1}}
\newtcolorbox{control}{colback=rojo!5,colframe=rojo,fonttitle=\bfseries,title={\faQuestionCircle~Preguntas de control}}
\newtcolorbox{codigo}[1][De la ecuación al código]{colback=verde!6,colframe=verde,fonttitle=\bfseries,title={\faCode~#1}}
\newtcolorbox{cierre}{colback=verde!6,colframe=verde,fonttitle=\bfseries,title={\faFlagCheckered~Cierre del módulo}}
\newtcolorbox{nota}[1]{colback=amarillo!12,colframe=amarillo!80!black,fonttitle=\bfseries,title={\faLightbulb~#1}}

% \cabecera{número}{título}{subtítulo}
\newcommand{\cabecera}[3]{%
  \title{\textcolor{morado}{\textbf{Ficha #1. #2}}\\[3pt]\large Taller virtual ATA EPE · #3}
  \author{}\date{}}

\cabecera{4}{Uniones que se sostienen: tornillos, tuercas y cuerpos compuestos}{Módulo 4: ligaduras dinámicas}

\begin{document}
\maketitle
\vspace{-2.5em}

\begin{meta}
\textbf{\textcolor{amarillo}{Meta del módulo.}} Hacer que las uniones se cumplan mientras el montaje se mueve:
calcular en cada paso los impulsos que las sostienen, sin que se separen con el tiempo. Modelar el tornillo como pieza,
los estados de la tuerca y la fusión de piezas en un solo cuerpo. Cara pública: \textbf{tornillos y tuercas} —
aflojar, ajustar y apretar cambian cómo se mueve el montaje.
\end{meta}

\begin{lectura}
\begin{enumerate}[leftmargin=*]
  \item \textbf{Baraff y Witkin, \emph{Physically Based Modeling} (SIGGRAPH 2001)}, notas de dinámica con
        ligaduras: la formulación de impulsos.
  \item \textbf{Catto, «Iterative dynamics with temporal coherence» (GDC 2005)}: impulsos secuenciales,
        acumulación y arranque en caliente. Es el corazón práctico del módulo.
  \item \textbf{Hadfield y Lasenby, «Constrained dynamics in conformal and projective geometric algebra» (2020)}:
        las ligaduras con potencia virtual en PGA.
  \item Macklin, Müller y Chentanez, «XPBD» (2016) y Müller et al., «Detailed rigid body simulation with XPBD»
        (2020): la alternativa en posiciones, para comparar.
  \item Goldstein, §2.4: multiplicadores de Lagrange (repaso del módulo 3).
\end{enumerate}
\end{lectura}

% =================================================================================================
\section{De la ecuación a los impulsos}

Con ligaduras, el movimiento cumple
\begin{equation}
  \mathbf M\dot{\vect u} = \vect f + \mathbf J^{\mathsf T}\vect\lambda,\qquad \mathbf J\vect u = \vect 0,
\end{equation}
donde $\vect u$ junta las velocidades de todos los cuerpos y $\mathbf M$ sus masas e inercias. Discretizando con el
paso semiimplícito del módulo 2, la velocidad nueva es $\vect u^+ = \vect u + h\mathbf M^{-1}(\vect f + \mathbf
J^{\mathsf T}\vect\lambda)$. Exigiendo $\mathbf J\vect u^+ = -\tfrac{\beta}{h}\vect C$ (cumplir la ligadura en
velocidad y corregir la deriva acumulada en posición):
\begin{equation}
  \bigl(\mathbf J\mathbf M^{-1}\mathbf J^{\mathsf T}\bigr)\,\vect P = -\mathbf J\bigl(\vect u + h\mathbf M^{-1}\vect f\bigr)
  - \frac{\beta}{h}\vect C,\qquad \vect P = h\vect\lambda .
\end{equation}
$\vect P$ son \emph{impulsos}. El término $\beta\vect C/h$ (estabilización de Baumgarte, $\beta\approx0{,}1$--$0{,}2$)
devuelve poco a poco las piezas a su lugar si se separaron por redondeo.

\begin{ejercicio}{4.1}
Para un péndulo (un punto de masa $m$ atado a distancia $\ell$ del origen), escribe $C$, $\mathbf J$ y la ecuación
(1.2) con un solo impulso. Verifica que el impulso es la tensión del hilo por $h$.
\end{ejercicio}

\subsection{Impulsos secuenciales (Gauss--Seidel proyectado)}
En lugar de resolver (1.2) entera, se recorren las ligaduras una por una, corrigiendo la velocidad con el impulso que
cumple esa sola fila:
\begin{equation}
  \Delta P_i = -\frac{\mathbf J_i\vect u + b_i}{\mathbf J_i\mathbf M^{-1}\mathbf J_i^{\mathsf T}},\qquad
  \vect u \leftarrow \vect u + \mathbf M^{-1}\mathbf J_i^{\mathsf T}\Delta P_i ,
\end{equation}
y se repite unas 8--20 veces. Es el método de Gauss--Seidel sobre (1.2). Se \emph{acumula} el impulso total de cada
fila; si la fila tiene límites (un motor con par máximo, la fricción), se recorta el acumulado, no el incremento. Al
empezar el paso siguiente se reutilizan los impulsos del anterior (\emph{arranque en caliente}): converge mucho más
rápido.

\begin{ejercicio}{4.2}
Muestra que (1.3) es exactamente una iteración de Gauss--Seidel para (1.2). ¿Por qué converge (pista: la matriz
$\mathbf J\mathbf M^{-1}\mathbf J^{\mathsf T}$ es simétrica y semidefinida positiva)? ¿Qué pasa si hay ligaduras
redundantes, como en el módulo 3?
\end{ejercicio}

% =================================================================================================
\section{Las uniones del kit en PGA}

Cada unión tiene un \emph{ancla} en cada pieza: la recta del agujero, $L_a$ y $L_b$, en sus marcos.
\begin{itemize}[leftmargin=*]
  \item \textbf{Posición}: las rectas deben coincidir en el mundo. El error $M_aL_a\rev M_a - M_bL_b\rev M_b$ es un
        bivector; sus componentes dan $\vect C$ (y para la fija se agrega el ángulo relativo).
  \item \textbf{Velocidad}: las filas de $\mathbf J$ son una base de los forques recíprocos a lo que la unión permite
        (ficha 3, §4). La unión de revolución da 5 filas; la cilíndrica, 4; la fija, 6.
  \item \textbf{Fricción en el eje}: una fila extra a lo largo de $L$ (velocidad relativa de giro igual a cero), con el
        impulso acumulado recortado a $\pm\tau_f\,h$. Si no hay par suficiente para vencerla, no gira.
\end{itemize}

% =================================================================================================
\section{Tornillos y tuercas}

\subsection{Estados}
\begin{tabularx}{\linewidth}{@{}l l Y@{}}
\toprule
\textbf{Estado de la tuerca} & \textbf{Unión pieza–tornillo} & \textbf{Comportamiento}\\
\midrule
Sin tuerca & cilíndrica & gira y se desliza; la gravedad puede sacar el tornillo\\
Floja & revolución & gira libre (fricción pequeña)\\
Ajustada & revolución con fricción $\tau_f$ & gira solo si se fuerza\\
Apretada & fija & las piezas se fusionan (§4)\\
\bottomrule
\end{tabularx}

\subsection{Apriete y enroscado}
Apretar con un par $T$ da una precarga $F_p = T/(Kd)$ y un par de fricción $\tau_f = \mu F_p r_m$ (plan, sección del
grafo de ensamble). Enroscar la tuerca es un movimiento de tornillo con el paso de la rosca: en PGA,
\begin{equation}
  M_{\text{tuerca}}(\theta) = \exp\!\Bigl(-\tfrac{\theta}{2}\bigl(L + \tfrac{p}{2\pi}\,\I L\bigr)\Bigr),
\end{equation}
con $L$ el eje del tornillo y $p = 0{,}7$\,mm. La animación del apriete es esta exponencial.

\begin{ejercicio}{4.3}
Verifica que (3.1) avanza $p$ por vuelta. ¿Qué cambia si la rosca fuera izquierda?
\end{ejercicio}

% =================================================================================================
\section{Cuerpos compuestos}

Las uniones fijas definen clases de equivalencia de piezas (unión–búsqueda). Cada clase es un solo cuerpo rígido, con
la masa e inercia combinadas (Steiner, módulo 2). Al fusionar, se conserva el momento total. En PGA es especialmente
limpio: los momentos son rectas en el espacio, y se suman.
\begin{equation}
  P_{\text{compuesto}} = \sum_i P_i^{(e)}\quad\Longrightarrow\quad \Omega_{\text{compuesto}} = \A_{\text{compuesto}}^{-1}[P_{\text{compuesto}}].
\end{equation}
Al aflojar, el cuerpo se separa y cada pieza hereda la velocidad rígida del cuerpo en su posición.

\begin{ejercicio}{4.4}
Muestra que al fusionar la energía cinética nunca aumenta: $T_{\text{después}}\le T_{\text{antes}}$ (es un choque
perfectamente inelástico). ¿Cuándo se conserva exactamente?
\end{ejercicio}

\begin{nota}{Pares de piezas que no chocan}
Las piezas unidas por un tornillo se tocan por diseño (el tornillo atraviesa el agujero). En el módulo 5, el detector
de colisiones debe ignorar los pares conectados en el grafo de ensamble.
\end{nota}

\begin{codigo}[Qué programar]
\begin{tabularx}{\linewidth}{@{}l Y@{}}
\texttt{nucleo/ligaduras.js} & Uniones de revolución, cilíndrica y fija; fila de fricción; solucionador por impulsos
  secuenciales con acumulación, recorte y arranque en caliente; Baumgarte.\\
\texttt{kit/ensamble.js} & Grafo de piezas y uniones; insertar un tornillo por agujeros alineados.\\
\texttt{kit/tornilleria.js} & Estados de la tuerca, apriete, enroscado (3.1).\\
\texttt{kit/compuestos.js} & Unión–búsqueda, fusión con conservación del momento (4.1), separación.\\
\texttt{app/} & Herramientas tornillo, tuerca y llave; contorno del cuerpo compuesto.\\
\end{tabularx}
\end{codigo}

\begin{control}
\begin{enumerate}[leftmargin=*]
  \item ¿Por qué se resuelven impulsos en lugar de fuerzas?
  \item ¿Qué hace la estabilización de Baumgarte y qué pasa si $\beta$ es muy grande?
  \item ¿Por qué se recorta el impulso acumulado y no el incremento?
  \item ¿Qué gana el arranque en caliente?
  \item ¿Por qué fusionar un tornillo con una rueda mejora la estabilidad numérica?
  \item ¿Qué se conserva al fusionar y qué no?
\end{enumerate}
\end{control}

\begin{cierre}
Ejercicios 4.1--4.4 · pruebas P2 (péndulo físico), P8 (fusión conserva el momento) y P9 (el tornillo sin tuerca se
cae) · cadena de 5 vigas estable · cara pública: tornillos y tuercas.
\end{cierre}

\end{document}
